\documentclass[11pt]{article}
\usepackage[margin=1in]{geometry}
\usepackage{amsmath,amssymb,amsthm}
\usepackage[T1]{fontenc}
\usepackage{lmodern}
\usepackage{xcolor}
\usepackage{listings}
\usepackage{booktabs}
\usepackage{longtable}
\usepackage[hidelinks]{hyperref}
\usepackage{xurl}

\newtheorem{theorem}{Theorem}
\newtheorem{lemma}[theorem]{Lemma}
\newtheorem{corollary}[theorem]{Corollary}
\newtheorem{proposition}[theorem]{Proposition}
\theoremstyle{remark}
\newtheorem{remark}[theorem]{Remark}

\newcommand{\lean}[1]{\texttt{#1}}
\DeclareUnicodeCharacter{2081}{\ensuremath{_1}}\DeclareUnicodeCharacter{2082}{\ensuremath{_2}}
\newcommand{\ip}[2]{\langle #1,#2\rangle}
\newcommand{\nm}[1]{\lVert #1\rVert}
\newcommand{\R}{\mathbb{R}}
\newcommand{\C}{\mathbb{C}}
\newcommand{\repoD}{https://github.com/jpbald93/dandelin}
\newcommand{\repoM}{https://github.com/jpbald93/marden}
% version-2 commits: dandelin=89309188e77f1564fd31cc9a2d245dbba1b34f8c marden=b4fbe75198e728699764a9b44e8f9981ba589299
\newcommand{\commitD}{205e4807a71d05898967d870e4528fb97d78b26d}
\newcommand{\commitM}{bd2cd3c9870905a04d38d079be4d83c824cbf9db}
\newcommand{\srcD}[2]{\href{\repoD/blob/\commitD/Dandelin/#1\#L#2}{\texttt{#1:#2}}}
\newcommand{\srcM}[2]{\href{\repoM/blob/\commitM/Marden/#1\#L#2}{\texttt{#1:#2}}}

\lstdefinelanguage{lean}{
  morekeywords={theorem,lemma,def,example,by,fun,have,show,exact,intro,namespace,end,import,let,noncomputable},
  sensitive=true, morecomment=[l]{--}, morecomment=[s]{/-}{-/}, morestring=[b]"}
\lstset{language=lean, basicstyle=\ttfamily\small, keywordstyle=\bfseries,
  columns=fullflexible, keepspaces=true, breaklines=true, frame=single, framerule=0.3pt,
  literate={ℝ}{{$\mathbb{R}$}}1 {ℂ}{{$\mathbb{C}$}}1 {≠}{{$\neq$}}1 {↔}{{$\leftrightarrow$}}2
           {→}{{$\to$}}1 {∀}{{$\forall$}}1 {∃}{{$\exists$}}1 {∧}{{$\wedge$}}1 {∨}{{$\vee$}}1
           {∈}{{$\in$}}1 {∉}{{$\notin$}}1 {≤}{{$\le$}}1 {•}{{$\bullet$}}1
           {⟪}{{$\langle\!\langle$}}2 {⟫}{{$\rangle\!\rangle$}}2 {‖}{{$\Vert$}}1
           {₀}{{$_0$}}1 {₁}{{$_1$}}1 {₂}{{$_2$}}1 {₃}{{$_3$}}1 {'}{{$'$}}1}

\title{Dandelin spheres and Marden's theorem,\\ formalised in Lean~4}
\author{Josh Bald\\ \small Independent researcher, Ontario, Canada\\ \small\texttt{jpbald93@gmail.com}}
\date{30 September 2026\\ \small Version~2, 22 October 2026: adds the Layer-14 and Layer-15 checks
(Section~\ref{sec:layers})}

\emergencystretch=3em
\begin{document}
\maketitle

\begin{abstract}
We formalise two classical theorems about conics in Lean~4 with Mathlib. The first is Dandelin's
theorem. A plane that misses the apex of a right circular cone meets it in an ellipse, a parabola
or a hyperbola. The spheres inscribed in the cone and tangent to the plane touch it at the foci.
In three-dimensional space we prove the classical picture: among the spheres centred on the axis
and touching the cone, exactly two are tangent to the plane in the ellipse and hyperbola cases and
exactly one in the parabola case; the section is \emph{exactly} the locus of constant focal sum
(ellipse) or constant absolute focal difference (hyperbola, on the double cone); for non-circular
sections the focus--directrix property holds with eccentricity $e=\cos\beta/\cos\alpha$, together
with its converse; the section is non-empty; the three regimes are characterised geometrically;
and the planes through the apex give a point, a line or two lines. The focal identities and
their converses hold in an arbitrary real inner product space; non-emptiness, the regime
characterisations and the two-line apex section are proved under explicit finite-dimension
assumptions. The second is Marden's theorem. For a non-degenerate triangle $a,b,c\in\C$, the ellipse
whose foci are the critical points of $(z-a)(z-b)(z-c)$ passes through the midpoints of the sides,
touches each side only there, and has the centroid as its centre. We also prove that it is the
only ellipse, given by foci and focal sum, with this property. Tangency is stated metrically, and
we explain why this is the right notion. The version-1 developments have about 2170 and 760 lines, respectively; version~2 adds 2287
Dandelin lines and 879 Marden lines in the Layer-14 and Layer-15 files. They
contain no \lean{sorry}, and their main theorems depend only on Lean's standard axioms. We found
no earlier complete formalisation of either theorem in the sources we searched;
Section~\ref{sec:related} records the scope of that search and one unresolved lead.
Version~2 restates the main theorems with Mathlib-only types, each with an explicit
non-degenerate witness, and compares the development with textbook statements formalised
independently; the comparison adds further proved statements and finds no false one.
\end{abstract}

\section{Introduction}

We work in Lean~4~\cite{Lean4} with Mathlib~\cite{MathlibCPP}. Pascal's hexagon theorem is
treated in a companion paper~\cite{BaldPascal}. This paper deals with
two other classical results about conics. Both are on the ``1000+ theorems'' list that Mathlib
maintains, and at the Mathlib version we use neither entry has a formal declaration
attached~\cite{Mathlib1000}.

\emph{Dandelin's theorem} (1822)~\cite{Dandelin} links the definition of a conic as a plane section of a cone to its
definition by foci. Suppose a plane cuts one nappe of a right circular cone in a closed curve.
There are two spheres that are inscribed in the cone and tangent to the plane, one on each side of it. If $F_1$, $F_2$ are the points of tangency,
then every point $X$ of the curve satisfies $|XF_1|+|XF_2|=\text{const}$. The proof uses only the fact
that all tangent segments from a point to a sphere have the same length. The same construction
gives the hyperbola, with one sphere in each nappe, and the parabola, with one sphere. The circle along which a sphere touches the cone lies
in a plane $\sigma$; for a non-circular section, $\sigma$ meets the cutting plane in the directrix. For textbook accounts
see Hilbert and Cohn-Vossen~\cite{HilbertCV}.

\emph{Marden's theorem}~\cite{Marden1945,MardenBook} is named after Marden, and it goes back to
Siebeck~\cite{Siebeck}. Kalman gives an elementary proof~\cite{Kalman}, and Minda and Phelps
give another~\cite{MindaPhelps}. It states
the following. Let $a,b,c\in\C$ be non-collinear and let $p(z)=(z-a)(z-b)(z-c)$. Then the roots of
$p'$ are the foci of an ellipse that is tangent to the three sides of the triangle at their
midpoints. This ellipse is the \emph{Steiner inellipse}, and its centre is the centroid.

\paragraph{Contributions.}
\begin{itemize}
\item A Lean formalisation of Dandelin's theorem for all three non-degenerate conics
(Section~\ref{sec:dandelin}). The focal identities hold in any real inner product space. The
hypotheses concern only the cone and the plane. The spheres, the foci and (for non-circular
sections) the directrices are constructed, and we prove that the tangent spheres are unique
within the family of axis spheres touching the cone (in $\R^3$, a Layer-15 lemma extends this to
every sphere that meets each generator in exactly one point; Section~\ref{sec:layers}). Every characterisation is proved in both directions: the section
\emph{equals} the focal locus.
\item The regime classification: the bounded-section inequality is equivalent to the section
lying in one nappe and being bounded, and the hyperbola inequality is equivalent to meeting both
nappes. We also prove non-emptiness, and treat the degenerate sections by planes through the apex.
These results assume that $E$ is finite-dimensional with the stated lower bound on its dimension,
except where noted.
\item A Lean formalisation of Marden's theorem with a metric notion of tangency, together with the
uniqueness of the inellipse among all ellipses given by two foci and a focal sum
(Section~\ref{sec:marden}).
\item Lemmas that help a reader check that the formal statements mean what they say
(Section~\ref{sec:correct}). They include a Lean counterexample showing that one of our
statements cannot be strengthened in the obvious way.
\item (Version~2.) Mathlib-only restatements of the main theorems with witnesses (Layer~14), and a
comparison with independently formalised textbook statements (Layer~15), Section~\ref{sec:layers}.
\end{itemize}

\paragraph{Availability.}
The code is at \url{\repoD} and \url{\repoM}. Line references are to commits
\texttt{205e480} and \texttt{bd2cd3c} respectively.
Version~2 adds the Layer-14 and Layer-15 files (\path{Dandelin/L14.lean},
\path{Dandelin/L15Classical.lean}, \path{Marden/L14.lean}, \path{Marden/L15Classical.lean}) in a
new commit of each repository (Dandelin \texttt{89309188e77f1564fd31cc9a2d245dbba1b34f8c}, Marden \texttt{b4fbe75198e728699764a9b44e8f9981ba589299}). That commit
also adds two import lines to each root file and updates \path{gate.sh} and the README. The
earlier Lean files and statements are unchanged, so the line references above remain valid.
The toolchain is Lean \texttt{v4.33.1}. Both projects pin Mathlib at tag \texttt{v4.33.1}
(revision \texttt{0df444a}) in \texttt{lake-manifest.json}.
Supporting material (the Layer-14 and Layer-15 reports, the blind writer's statements and report, the transcript audit, and the reviews of this version with the fixes made in response) is in the folder \texttt{Conics-dandelin-marden} of the repository \url{https://github.com/jpbald93/LEAN-verification}.

\begin{longtable}{@{}p{0.42\textwidth}p{0.29\textwidth}l@{}}
\caption{The main formal statements. Namespaces \lean{Dandelin} (upper part) and \lean{Marden}
(lower part). ``Thm'' refers to the numbered results of this paper.}\label{tab:thms}\\
\toprule
Lean name & In this paper & Source \\ \midrule
\endfirsthead
\toprule
Lean name & In this paper & Source \\ \midrule
\endhead
\bottomrule
\endfoot
\lean{tangent\_lengths\_eq} & equal tangents, \S\ref{sec:cone} & \srcD{Basic.lean}{29} \\
\lean{cone\_tangent\_sq} & Lemma~\ref{lem:conetan} & \srcD{Basic.lean}{53} \\
\lean{cone\_outside\_sphere} & Lemma~\ref{lem:inscribed} & \srcD{Basic.lean}{63} \\
\lean{tangencyPoint\_mem\_sphere} & Lemma~\ref{lem:inscribed} & \srcD{Basic.lean}{85} \\
\lean{tangencyPoint\_perp} & Lemma~\ref{lem:inscribed} & \srcD{Basic.lean}{95} \\
\lean{dandelin\_of\_between\_isTangentAt} & Remark~\ref{rem:mathlibtangent} & \srcD{Basic.lean}{146} \\
\lean{dandelin} & Theorem~\ref{thm:ellipse}(iv), $\Rightarrow$, normalised & \srcD{Basic.lean}{216} \\
\lean{tangent\_iff} & Theorem~\ref{thm:ellipse}(i), normalised ($s<m$, $c>0$) & \srcD{Existence.lean}{64} \\
\lean{dandelin\_converse} & Theorem~\ref{thm:ellipse}(iv), $\Leftarrow$ & \srcD{Existence.lean}{155} \\
\lean{dandelin\_exists} & Theorem~\ref{thm:ellipse} & \srcD{Existence.lean}{230} \\
\lean{section\_bounded} & Theorem~\ref{thm:regime} & \srcD{Existence.lean}{277} \\
\lean{infDist\_directrix} & Lemma~\ref{lem:infdist} & \srcD{Directrix.lean}{72} \\
\lean{focal\_dist\_eq\_ecc\_mul} & Theorem~\ref{thm:directrix} & \srcD{Directrix.lean}{148} \\
\lean{ecc\_lt\_one\_iff} & Theorem~\ref{thm:directrix}, $e<1$ & \srcD{Directrix.lean}{160} \\
\lean{ecc\_eq\_one} & Theorem~\ref{thm:directrix}, $e=1$ & \srcD{Directrix.lean}{168} \\
\lean{one\_lt\_ecc\_iff} & Theorem~\ref{thm:directrix}, $e>1$ & \srcD{Directrix.lean}{174} \\
\lean{tangent\_of\_proj} & Lemma~\ref{lem:proj} & \srcD{Directrix.lean}{192} \\
\lean{directrix\_ellipse} & Theorem~\ref{thm:directrix} & \srcD{Directrix.lean}{206} \\
\lean{directrix\_parabola} & Theorem~\ref{thm:parabola} & \srcD{Directrix.lean}{247} \\
\lean{mem\_cone2\_iff} & \S\ref{sec:conv} & \srcD{Hyperbola.lean}{25} \\
\lean{hyp\_upper} & Theorem~\ref{thm:hyperbola}(ii), upper & \srcD{Hyperbola.lean}{136} \\
\lean{hyp\_lower} & Theorem~\ref{thm:hyperbola}(ii), lower & \srcD{Hyperbola.lean}{157} \\
\lean{hyp\_converse} & Theorem~\ref{thm:hyperbola}(iii), $\Leftarrow$ & \srcD{Hyperbola.lean}{192} \\
\lean{dandelin\_hyperbola} & Theorem~\ref{thm:hyperbola} & \srcD{Hyperbola.lean}{283} \\
\lean{hyp\_both\_nappes} & Theorem~\ref{thm:regime}(c) & \srcD{Regime.lean}{222} \\
\lean{ellipse\_nonempty} & Theorem~\ref{thm:regime}(d) & \srcD{Regime.lean}{255} \\
\lean{section\_unbounded} & Theorem~\ref{thm:regime}(b) & \srcD{Regime.lean}{330} \\
\lean{one\_nappe} & Theorem~\ref{thm:regime}(a) & \srcD{Regime.lean}{373} \\
\lean{bounded\_iff} & Theorem~\ref{thm:regime}(b) & \srcD{Regime.lean}{407} \\
\lean{ellipse\_regime\_iff} & Theorem~\ref{thm:regime}(b) & \srcD{Regime.lean}{426} \\
\lean{hyperbola\_regime\_iff} & Theorem~\ref{thm:regime}(c) & \srcD{Regime.lean}{442} \\
\lean{parabola\_regime} & Theorem~\ref{thm:regime}(e) & \srcD{Regime.lean}{451} \\
\lean{regime\_R3} & Theorem~\ref{thm:regime} in $\R^3$ & \srcD{Regime.lean}{475} \\
\lean{directrix\_converse} & Theorem~\ref{thm:dirconv} & \srcD{Degenerate.lean}{30} \\
\lean{directrix\_iff\_ellipse} & Theorem~\ref{thm:dirconv}(a) & \srcD{Degenerate.lean}{88} \\
\lean{directrix\_iff\_parabola} & Theorem~\ref{thm:dirconv}(a) & \srcD{Degenerate.lean}{113} \\
\lean{directrix\_iff\_hyperbola} & Theorem~\ref{thm:dirconv}(b) & \srcD{Degenerate.lean}{133} \\
\lean{apex\_section\_point} & Theorem~\ref{thm:apex}(a) & \srcD{Degenerate.lean}{226} \\
\lean{apex\_section\_line} & Theorem~\ref{thm:apex}(b) & \srcD{Degenerate.lean}{252} \\
\lean{apex\_section\_two\_lines} & Theorem~\ref{thm:apex}(c) & \srcD{Degenerate.lean}{369} \\
\lean{apex\_two\_lines\_subset} & Theorem~\ref{thm:apex}, remark & \srcD{Degenerate.lean}{443} \\
\lean{apex\_section\_R3} & Theorem~\ref{thm:apex} in $\R^3$ & \srcD{Degenerate.lean}{470} \\
\midrule
\lean{midpoint\_mul} & Lemma~\ref{lem:key} & \srcM{Basic.lean}{24} \\
\lean{focalSum\_midpoint\_sq} & Lemma~\ref{lem:focsq} & \srcM{Basic.lean}{37} \\
\lean{focalSum\_midpoints\_eq} & Theorem~\ref{thm:marden}(i) & \srcM{Basic.lean}{59} \\
\lean{focal\_le\_of\_mul} & Lemma~\ref{lem:reflect} & \srcM{Tangent.lean}{73} \\
\lean{focal\_lt\_of\_mul} & Lemma~\ref{lem:reflect}, strict & \srcM{Tangent.lean}{100} \\
\lean{vieta} & \S\ref{sec:mardenstmt} & \srcM{Main.lean}{31} \\
\lean{exists\_critical\_points} & Lemma~\ref{lem:mardennonvac} & \srcM{Main.lean}{44} \\
\lean{marden} & Theorem~\ref{thm:marden} & \srcM{Main.lean}{136} \\
\lean{reflection\_of\_min} & Lemma~\ref{lem:convtan} & \srcM{Unique.lean}{74} \\
\lean{foci\_of\_reflection} & Lemma~\ref{lem:foci} & \srcM{Unique.lean}{194} \\
\lean{steiner\_inellipse\_unique} & Theorem~\ref{thm:unique}, weak form & \srcM{Unique.lean}{233} \\
\lean{nondeg\_of\_midpoints} & Lemma~\ref{lem:mardennonvac} & \srcM{Unique.lean}{322} \\
\lean{steiner\_unique} & Theorem~\ref{thm:unique} & \srcM{Unique.lean}{354} \\
\lean{steiner\_inellipse\_iff} & Corollary~\ref{cor:iff} & \srcM{Unique.lean}{373} \\
\end{longtable}

\section{Conventions}\label{sec:conv}

\paragraph{Ambient space.} Throughout the Dandelin development, $E$ is a real inner product space
(\lean{[NormedAddCommGroup E] [InnerProductSpace $\R$ E]}), and points are vectors of~$E$. There
is no finite-dimensionality assumption unless one is stated. Where one is stated, it is
\lean{[FiniteDimensional $\R$ E]} together with a bound on \lean{Module.finrank $\R$ E}, and we write
``$E$ finite-dimensional, $\dim E\ge k$''. The intended case is
\lean{EuclideanSpace $\R$ (Fin 3)}, and results that need a dimension bound are also specialised to
it.

\paragraph{Cone.} A cone is given by an apex $V$, a unit axis vector~$u$, and $k=\cos\alpha$,
$s=\sin\alpha$, where $\alpha$ is the half-angle. The hypotheses are $k^2+s^2=1$, $k>0$ and
(usually) $s>0$, so that $0<\alpha<\pi/2$. One nappe and the double cone are
\[
\lean{cone V u k}=\{X : \ip{X-V}{u}=k\nm{X-V}\},\qquad
\lean{cone2 V u k}=\{X : \ip{X-V}{u}^2=k^2\nm{X-V}^2\}.
\]
Both contain~$V$. \lean{mem\_cone2\_iff} proves that $\lean{cone2 V u k}=\lean{cone V u k}\cup\lean{cone V (-u) k}$.
The axis sphere with parameter $d\in\R$ has centre $C_d=V+du$ and radius $|d|s$, which is the
distance from $C_d$ to every generator. For $d>0$ it is inscribed in the upper nappe, and for
$d<0$ in the lower one (Lemma~\ref{lem:inscribed}). All uniqueness statements about spheres in
the main files are about this one-parameter family $\{(C_d,|d|s)\}$; the main files do not classify arbitrary
centres and radii satisfying an independent notion of being inscribed. The Layer-15 file adds this
classification in $\R^3$: a sphere of positive radius that meets every generator of the double
cone in exactly one point is an axis sphere $(C_d,|d|s)$ with $d\ne0$
(\lean{Dandelin.Classical.l15\_tangent\_cone\_iff}, Section~\ref{sec:layers}).

\paragraph{Plane.} $\lean{plane p}_0\lean{ n}=\{X:\ip{X-p_0}{n}=0\}$ with $\nm n=1$. We write
\[
m=\ip un,\qquad c=\ip{p_0-V}n .
\]
So $|m|=\sin\beta$, where $\beta$ is the angle between the plane and the axis, and $c$ is the signed
distance from $V$ to the plane. The plane misses the apex exactly when $c\neq0$.

\paragraph{Tangency.} A sphere with centre $C$ and radius $r$ is tangent to the plane when
$|\ip{C-p_0}n|=r$, that is, when the distance from the centre to the plane equals the radius. Its tangency point is then the projection
$F=C-\ip{C-p_0}n\,n$ of the centre onto the plane. Lemma~\ref{lem:proj} shows that this
matches the pointwise definition used elsewhere: $F$ is on the sphere, and $X-F\perp F-C$ for
every $X$ in the plane.

\paragraph{Distances to sets.} Distances to the directrix use Mathlib's \lean{Metric.infDist}.
Because $\lean{infDist}\,X\,\emptyset=0$, a statement of the form $|XF|=e\cdot\lean{infDist}\,X\,\ell$ could hold
vacuously if $\ell$ were empty. Lemma~\ref{lem:infdist} rules this out, since it computes the value
explicitly.

\paragraph{Marden.} Points are complex numbers. $\lean{focalSum f}_1\lean{ f}_2\lean{ z}=|z-f_1|+|z-f_2|$, and the ellipse with foci
$f_1,f_2$ and major axis $2A$ is $\{z:\lean{focalSum}\ f_1\,f_2\,z=2A\}$.

\section{Dandelin's theorem}\label{sec:dandelin}

\subsection{The cone model}\label{sec:cone}

The classical proof uses one fact about spheres: two tangent segments from a point~$X$ to a
sphere have the same length, since both have squared length $|XC|^2-r^2$ by Pythagoras
(\lean{tangent\_lengths\_eq}). The second fact is specific to the cone.

\begin{lemma}\label{lem:conetan}
If $X\in\lean{cone V u k}$, then $\nm{X-C_d}^2-(ds)^2=(\nm{X-V}-dk)^2$.
\end{lemma}

So, for $d>0$ and a non-apex point $X$ of the upper nappe, the tangent length from $X$ to the axis sphere with parameter~$d$ is
$|\nm{X-V}-dk|$, the distance along the generator from $X$ to the circle at distance $dk$ from the
apex. This is the locus of contact: the intersection of the sphere with the hyperplane $\sigma_d$ of
Section~\ref{sec:directrix}. When $E$ is finite-dimensional of dimension $N\ge2$ and $d\ne0$, it is a
sphere of dimension $N-2$ with centre $V+dk^2u$ and radius $|d|sk$; for $N=3$ it is a circle. The
lower nappe is symmetric. We call $\lean{tangencyPoint V X k d}=V+\frac{dk}{\nm{X-V}}(X-V)$ the point where the generator through $X$ touches the sphere.
This point lies on the sphere, and the generator is perpendicular to the radius there
(Lemma~\ref{lem:inscribed}). As a consequence, if a plane through~$X$ touches the sphere at~$F$, then
$|XF|=|\nm{X-V}-dk|$ (\lean{dist\_tangent\_eq}).

\subsection{The ellipse: existence, uniqueness and both directions}

\begin{theorem}[Dandelin, ellipse]\label{thm:ellipse}
Let $s<|m|$ and $mc>0$. Then there are $0<d_1<d_2$ and points $F_1,F_2$ such that:
\begin{enumerate}
\item[(i)] for every $d\in\R$, the sphere $(C_d,ds)$ is tangent to the plane if and only if
$d\in\{d_1,d_2\}$;
\item[(ii)] the two centres lie on opposite sides of the plane;
\item[(iii)] $F_i$ lies in the plane and is the projection of $C_{d_i}$ onto it;
\item[(iv)] for every $X$ in the plane, $X\in\lean{cone V u k}$ if and only if
$|XF_1|+|XF_2|=(d_2-d_1)k$.
\end{enumerate}
\end{theorem}

\begin{lstlisting}
theorem dandelin_exists {V u n p₀ : E} {k s : ℝ} (hu : ‖u‖ = 1) (hn : ‖n‖ = 1)
    (hks : k ^ 2 + s ^ 2 = 1) (hk : 0 < k) (hs : 0 < s) (hm : s < |⟪u, n⟫|)
    (hc : 0 < ⟪u, n⟫ * ⟪p₀ - V, n⟫) :
    ∃ d₁ d₂ : ℝ, ∃ F₁ F₂ : E, 0 < d₁ ∧ d₁ < d₂ ∧
      (∀ d : ℝ, |⟪(V + d • u) - p₀, n⟫| = d * s ↔ d = d₁ ∨ d = d₂) ∧
      ⟪(V + d₁ • u) - p₀, n⟫ * ⟪(V + d₂ • u) - p₀, n⟫ < 0 ∧
      F₁ ∈ plane p₀ n ∧ F₂ ∈ plane p₀ n ∧
      F₁ = (V + d₁ • u) - ⟪(V + d₁ • u) - p₀, n⟫ • n ∧
      F₂ = (V + d₂ • u) - ⟪(V + d₂ • u) - p₀, n⟫ • n ∧
      ∀ X ∈ plane p₀ n, (X ∈ cone V u k ↔ ‖X - F₁‖ + ‖X - F₂‖ = (d₂ - d₁) * k)
\end{lstlisting}

The two hypotheses concern only the cone and the plane. The first, $s<|m|$, says that the plane
is less steep than a generator. The second, $mc>0$, says that the plane misses the apex and is
on the side of the nappe \lean{cone V u k}. Theorem~\ref{thm:regime} shows that together they
characterise the ellipse regime.

\begin{proof}[Proof sketch]
Flip $n$ if necessary so that $m>0$ and $c>0$. The tangency condition $|c-dm|=ds$ is piecewise linear
in~$d$, and its solutions are $d_1=c/(m+s)$ and $d_2=c/(m-s)$ (\lean{tangent\_iff}). Take
$F_i$ to be the projections of the centres. For the forward direction of~(iv), let $X$ be a point
of the section and write $t=\nm{X-V}$. By Lemma~\ref{lem:conetan}, $|XF_1|+|XF_2|=|t-d_1k|+|t-d_2k|$.
This equals $(d_2-d_1)k$ if and only if $d_1k\le t\le d_2k$. Classically this ``betweenness'' is read off a
picture. Here it is proved (\lean{between\_of\_gram}). Apply the Cauchy--Schwarz inequality to the
components of $X-V$ and of $n$ orthogonal to~$u$. Combined with the tangency equations
$\ip{X-V}n=d_1(m+s)=d_2(m-s)$, this gives $d_1k\le t\le d_2k$.
For the converse, let $A=|XF_1|$, $B=|XF_2|$ and $p=\ip{X-V}u$. For $X$ in the plane,
Pythagoras gives $A^2=t^2-2d_1p+d_1^2k^2$, and similarly for $B$ with~$d_2$. Subtract the two
equations and use $A+B=(d_2-d_1)k$. This gives $kA=p-d_1k^2$, and hence $p>0$ and
$(p-kt)(p+kt)=0$. So $p=kt$, which is the cone equation (\lean{converse\_alg},
\lean{dandelin\_converse}).
\end{proof}

\begin{remark}[Methods]\label{rem:cos}
An earlier draft of the specification stated the bounded-section hypothesis as
$|\ip nu|>\cos\alpha$. Before the existence theorem was formalised, a numerical experiment on 2000 random
cone--plane pairs (\texttt{code/check\_phase2.py}) compared the prediction ``every generator of
the nappe meets the plane'' with the candidate inequality. It showed that the correct condition is
$|\ip nu|>\sin\alpha$. A generator direction $w=ku+se$ with $e\perp u$ a unit vector has
$\ip wn$ ranging over $[km-s\sqrt{1-m^2},\,km+s\sqrt{1-m^2}]$, and this interval avoids $0$
exactly when $m^2>s^2$. With the $\cos\alpha$ hypothesis, the existence statement is false for
$\alpha>\pi/4$, where it admits parabolic and hyperbolic sections. A failed proof attempt would
not have revealed this as directly.
\end{remark}

\subsection{The hyperbola}

In the hyperbola regime the plane cuts both nappes, so the statement has to be about the double
cone. Take one sphere in each nappe, that is, axis spheres with parameters of opposite signs.

\begin{theorem}[Dandelin, hyperbola]\label{thm:hyperbola}
Let $|m|<s$ and $c\neq0$. Then $e:=\sqrt{1-m^2}/k>1$, and there are $d_2<0<d_1$ and points
$F_1,F_2$ such that:
\begin{enumerate}
\item[(i)] for every $d\in\R$, the sphere $(C_d,|d|s)$ is tangent to the plane if and only if $d\in\{d_1,d_2\}$, and
$F_i$ is the projection of $C_{d_i}$ onto the plane;
\item[(ii)] for every $X$ in the cutting plane, $|XF_2|-|XF_1|=D$ if $X$ is on the upper nappe, and
$|XF_1|-|XF_2|=D$ if $X$ is on the lower nappe, where $D=(d_1-d_2)k$;
\item[(iii)] for every $X$ in the plane, $X\in\lean{cone2 V u k}$ if and only if $\bigl||XF_1|-|XF_2|\bigr|=D$;
\item[(iv)] every $X$ of the section satisfies $|XF_i|=e\cdot\mathrm{dist}(X,\ell_i)$ for
$i=1,2$, where $\ell_i$ is the directrix of Section~\ref{sec:directrix}.
\end{enumerate}
\end{theorem}

This is \lean{dandelin\_hyperbola}. Its orientation step flips $n$ so that $c>0$, and then
$d_1=c/(m+s)$ and $d_2=c/(m-s)$ as before. The closed form is
$d_1-d_2=2sc/(s^2-m^2)$ (\lean{hyp\_param\_sub}). The proof of~(ii) needs the analogue of
betweenness: a point of the upper nappe in the plane lies beyond the upper circle of contact,
$d_1k\le\nm{X-V}$ (\lean{hyp\_between}). The proof again uses Cauchy--Schwarz. The lower-nappe identity follows from
the upper one under the symmetry $u\mapsto-u$, $d_i\mapsto-d_{3-i}$.

For the converse (\lean{hyp\_converse}), write $\bigl||XF_1|-|XF_2|\bigr|=D$ as
$|XF_1|-|XF_2|=\varepsilon D$ with $\varepsilon=\pm1$. The same subtraction as in the ellipse case
gives $\varepsilon kA=d_1k^2-p$, and then $p^2=k^2t^2$, which is the equation of \lean{cone2}. The
sign~$\varepsilon$ tells us which nappe $X$ lies on, so no positivity argument is needed. The
converse cannot be stated for one nappe. Each branch of the hyperbola lies on its own nappe, and the
absolute value is necessary.

\subsection{Focus and directrix; the parabola}\label{sec:directrix}

Let $\sigma_d=\{Y:\ip{Y-V}u=dk^2\}$ be the plane perpendicular to the axis at height~$dk^2$.
The circle of contact of the sphere with parameter~$d$ lies in~$\sigma_d$
(\lean{tangencyPoint\_mem\_axisPlane}). The \emph{directrix} is
$\ell_d=\lean{directrix V u n p}_0\,(dk^2)=\lean{plane p}_0\lean{ n}\cap\sigma_d$. The eccentricity is
\[
e=\lean{ecc}\ m\ k=\frac{\sqrt{1-m^2}}{k}=\frac{\cos\beta}{\cos\alpha}.
\]

\begin{lemma}\label{lem:infdist}
If $m^2<1$ and $X$ lies in the plane, then
$\lean{infDist}\,X\,\ell_d=|\ip{X-V}u-dk^2|/\sqrt{1-m^2}$.
\end{lemma}

The proof exhibits the nearest point $X-\frac{\delta}{1-m^2}(u-mn)$ explicitly, where
$\delta=\ip{X-V}u-dk^2$, and uses Cauchy--Schwarz against $u-mn$ for the lower bound. It holds
in every inner product space. In particular $\ell_d$ is non-empty whenever the plane is.

\begin{theorem}[focus--directrix]\label{thm:directrix}
Suppose $m^2<1$, $X$ lies in the plane and on \lean{cone V u k}, and a sphere $(C_d,ds)$ touches a
plane through $X$ at~$F$. Then $|XF|=e\cdot\lean{infDist}\,X\,\ell_d$. Moreover $e<1$ iff $s<|m|$,
$e=1$ if $s=|m|$, and $e>1$ iff $|m|<s$. In the ellipse regime (the hypotheses of
Theorem~\ref{thm:ellipse} together with $m^2<1$) both foci satisfy this relation with their own
directrices.
\end{theorem}

The proof composes two elementary facts. The first is $|XF|=\mathrm{dist}(X,\sigma_d)/\cos\alpha$
(\lean{focal\_dist\_eq\_infDist\_axisPlane}): the tangent length along the generator is $1/\cos\alpha$
times its axial component. The second is
$\mathrm{dist}(X,\sigma_d)=\cos\beta\cdot\mathrm{dist}(X,\ell_d)$ for $X$ in the plane
(\lean{infDist\_axisPlane\_eq\_mul}). The hypothesis $m^2<1$ excludes the plane perpendicular to
the axis. There the section is a circle and $\ell_d$ is empty.

\begin{theorem}[parabola]\label{thm:parabola}
Let $s=|m|$ and $mc>0$. Then $e=1$, and there is exactly one $d>0$ for which $(C_d,ds)$ is tangent to the plane.
Its tangency point~$F$ satisfies $|XF|=\lean{infDist}\,X\,\ell_d$ for every $X$ in the section
(\lean{directrix\_parabola}).
\end{theorem}

The converse of the focus--directrix property also holds, and it turns out to be a statement
about the double cone. Let $F$ be the tangency point of $(C_d,|d|s)$ with the plane and let $X$ be
any point of the plane. Combining Lemma~\ref{lem:infdist} with Pythagoras (\lean{focal\_sq}) gives
\[
(k|XF|)^2-\bigl(\ip{X-V}u-dk^2\bigr)^2=k^2\nm{X-V}^2-\ip{X-V}u^2 .
\]
So the relation $|XF|=e\,\mathrm{dist}(X,\ell_d)$ restricted to the plane \emph{is} the equation
of \lean{cone2}.

\begin{theorem}[converse]\label{thm:dirconv}
If $m^2<1$, $X$ is in the plane, the plane touches $(C_d,|d|s)$ at $F$ and
$|XF|=e\cdot\lean{infDist}\,X\,\ell_d$, then $X\in\lean{cone2 V u k}$
(\lean{directrix\_converse}; the formal lemma needs only the pointwise conditions $|FC_d|=|d|s$ and
$\ip{X-F}{F-C_d}=0$, not tangency of the whole plane, and does not assume $s>0$). Moreover:
\begin{enumerate}
\item[(a)] in the ellipse and parabola regimes, the section equals the one-nappe locus: $X\in\lean{cone V u k}$ if and
only if the relation holds, for each focus (\lean{directrix\_iff\_ellipse},
\lean{directrix\_iff\_parabola});
\item[(b)] in the hyperbola regime, $X\in\lean{cone2 V u k}$ if and only if the relation holds, for each
focus. The one-nappe version is false: there is a point $Y$ of the plane that satisfies the relation for
\emph{both} foci but is not on \lean{cone V u k}, and there is another that is not on \lean{cone V (-u) k}
(\lean{directrix\_iff\_hyperbola}).
\end{enumerate}
\end{theorem}

The counterexamples in~(b) are proved in Lean, not merely observed. The witnesses are the section
points on the other nappe constructed in Theorem~\ref{thm:regime}(c). They are not the apex,
because $c\ne0$ and a point lying on both nappes must be the apex (\lean{eq\_of\_mem\_both\_nappes}).

\subsection{Regimes and non-emptiness}\label{sec:regime}

So far the conditions $s<|m|$, $s=|m|$ and $|m|<s$ have only been hypotheses. The following
theorem shows that they have the intended geometric meaning. All of its parts assume the standing
hypotheses on $u,n,k,s$ and $c\ne0$.

\begin{theorem}[regimes]\label{thm:regime}
\leavevmode
\begin{enumerate}
\item[(a)] If $s\le|m|$, the section of the double cone lies in a single nappe: in
\lean{cone V u k} if $mc>0$ and in \lean{cone V (-u) k} if $mc<0$ (\lean{one\_nappe}; any~$E$).
\item[(b)] If $E$ is finite-dimensional and $\dim E\ge3$, then $s<|m|$ holds iff $\lean{plane}\cap\lean{cone2}$ is bounded
(\lean{bounded\_iff}), iff it lies in one nappe and is bounded (\lean{ellipse\_regime\_iff}).
The forward direction holds in any~$E$ (\lean{ellipse\_regime}, via \lean{section\_bounded}).
\item[(c)] $|m|<s$ holds iff the plane meets both nappes, \lean{cone V u k} and
\lean{cone V (-u) k} (\lean{hyperbola\_regime\_iff}; any~$E$).
\item[(d)] If $E$ is finite-dimensional, $\dim E\ge2$ and $s<|m|$, $mc>0$, then the section is non-empty (\lean{ellipse\_nonempty}). The witness is
the intersection of the plane with a generator of the nappe.
\item[(e)] If $E$ is finite-dimensional, $\dim E\ge3$ and $|m|=s$, the section lies in one nappe and is unbounded
(\lean{parabola\_regime}), and in particular it is non-empty.
\end{enumerate}
\end{theorem}

The proofs give explicit geometric witnesses. (This is not a claim about computational content:
like most of Mathlib, the developments use \lean{Classical.choice}.) A generator $V+t(ku+se)$ with $e\perp u$ a unit vector meets the plane at
$t=c/\ip{ku+se}n$, and the sign of $t$ decides the nappe (\lean{gen\_point}). For the
hyperbola, the directions $e=\pm q/\nm q$ with $q=n-mu$ give parameters of opposite signs. For
unboundedness, choose a unit $f\perp u,n$ (this is where finite dimension and $\dim E\ge3$ are used,
\lean{exists\_unit\_orth}) and tilt $e$ in the $(q,f)$-plane so that $\ip{ku+se}n$ becomes an
arbitrarily small nonzero number (\lean{section\_unbounded}).

The dimension hypotheses are necessary. If $\dim E=1$, the cone is $\{V\}$, so the section is
empty. If $\dim E=2$, a ``plane'' is a line and the double cone is a pair of lines through~$V$, so
every section is bounded, even when $|m|\le s$. These two remarks are elementary and are not
formalised. \lean{regime\_R3} bundles (b), (c) and~(d) for \lean{EuclideanSpace $\R$ (Fin 3)}.

\subsection{Planes through the apex}

\begin{theorem}\label{thm:apex}
Let $c=0$ and put $a=u-mn$ (\lean{apexDir u n}).
\begin{enumerate}
\item[(a)] If $s<|m|$, then $\lean{plane}\cap\lean{cone2}=\{V\}$ (any $E$).
\item[(b)] If $|m|=s$, then $a\ne0$ and the section is the line $\{V+ta\}$ (any $E$).
\item[(c)] If $|m|<s$ and $\dim E=3$, then there are a unit vector $f\perp u,n$ and $\rho>0$ such that the section is
the union of the two lines $\{V+t(a\pm\rho f)\}$, and these lines meet only at~$V$.
\end{enumerate}
\end{theorem}

Here $\rho=\sqrt{(1-m^2)(s^2-m^2)}/k$ (\lean{apexRho}). In case~(b), equality in Cauchy--Schwarz
forces $X-V$ to be parallel to~$a$. In case~(c), $\dim E=3$ is used to decompose every
$w\perp n$ in the orthogonal pair $a,f$ (\lean{decomp\_of\_finrank\_le\_three}). Four nonzero
pairwise orthogonal vectors would be linearly independent. Then the cone equation factors as
$k^2(\beta-\alpha\rho)(\beta+\alpha\rho)=0$, where $\alpha$ and $\beta$ are the coordinates of $w$
along $a$ and~$f$. If $E$ is finite-dimensional with $N=\dim E>3$, the section is not two lines but
a quadric cone in the $(N-1)$-dimensional cutting hyperplane, of locus dimension $N-2$. (For
example, in $\R^4$ with $u=e_4$, $n=e_1$ and $k=s=1/\sqrt2$, the section is
$\{x_1=0,\ x_4^2=x_2^2+x_3^2\}$, a two-dimensional cone in a three-dimensional hyperplane.) This
description is informal. Only the containment of the two lines in the section is formalised
(\lean{apex\_two\_lines\_subset}, for $E$ finite-dimensional with $\dim E\ge3$).

\section{Marden's theorem}\label{sec:marden}

\subsection{Statement}\label{sec:mardenstmt}

Let $\lean{cubic a b c}=(X-a)(X-b)(X-c)\in\C[X]$. We state the critical points through Mathlib's
formal derivative. The hypothesis is
\[
\lean{derivative (cubic a b c)} = 3(X-z_1)(X-z_2),
\]
that is, $z_1,z_2$ are the roots of $p'$ counted with multiplicity. \lean{vieta} derives
$z_1+z_2=\tfrac23(a+b+c)$ and $z_1z_2=\tfrac13(ab+bc+ca)$ from it. Non-collinearity is stated as
$\operatorname{Im}\frac{c-a}{b-a}\ne0$.

\begin{theorem}[Marden]\label{thm:marden}
Let $a,b,c$ be distinct and non-collinear, and let $z_1,z_2$ be as above. Put
$2A=\lean{focalSum}\ z_1\,z_2\,\frac{a+b}2$. Then:
\begin{enumerate}
\item[(i)] all three midpoints lie on the ellipse $\{z:|z-z_1|+|z-z_2|=2A\}$;
\item[(ii)] for each side $xy$ and every real $t\ne0$,
$2A<\lean{focalSum}\ z_1\,z_2\,\bigl(\frac{x+y}2+t(y-x)\bigr)$;
\item[(iii)] $\frac{z_1+z_2}2=\frac{a+b+c}3$.
\end{enumerate}
\end{theorem}

\begin{lstlisting}
theorem marden {a b c z₁ z₂ : ℂ}
    (hder : derivative (cubic a b c) = C 3 * (X - C z₁) * (X - C z₂))
    (hab : a ≠ b) (hbc : b ≠ c) (hca : c ≠ a)
    (hnc : ((c - a) / (b - a)).im ≠ 0) :
    let A := focalSum z₁ z₂ ((a + b) / 2) / 2
    (focalSum z₁ z₂ ((a + b) / 2) = 2 * A ∧
      focalSum z₁ z₂ ((b + c) / 2) = 2 * A ∧
      focalSum z₁ z₂ ((c + a) / 2) = 2 * A) ∧
    (∀ t : ℝ, t ≠ 0 → 2 * A < focalSum z₁ z₂ ((a + b) / 2 + t * (b - a))) ∧
    (∀ t : ℝ, t ≠ 0 → 2 * A < focalSum z₁ z₂ ((b + c) / 2 + t * (c - b))) ∧
    (∀ t : ℝ, t ≠ 0 → 2 * A < focalSum z₁ z₂ ((c + a) / 2 + t * (a - c))) ∧
    (z₁ + z₂) / 2 = (a + b + c) / 3
\end{lstlisting}

The distinctness hypotheses follow from \lean{hnc}. We keep them for readability.

\subsection{Why metric tangency is tangency}\label{sec:metrictan}

Condition~(ii) says that the side line through the midpoint $m$ meets the closed region
$K=\{z:\lean{focalSum}\ z_1\,z_2\,z\le2A\}$ only at~$m$. The ellipse is non-degenerate:
$2A>|z_1-z_2|$ by Lemma~\ref{lem:mardennonvac}. So $K$ is a compact, strictly convex set with
non-empty interior, and its boundary is a smooth curve. For such a set, a line through a boundary
point~$m$ that is not tangent at~$m$ enters the interior of~$K$ on one side of~$m$, so it contains
points with focal sum $<2A$. The tangent line at~$m$ supports~$K$ and, by strict convexity, meets
$K$ only at~$m$. Hence~(ii) is equivalent to ``the side line is tangent to the ellipse at the
midpoint''. For \emph{inscribed}, argue as follows. By~(ii), $K$ lies in one of the two closed
half-planes bounded by each side line. The midpoint of another side lies in~$K$ and strictly
inside the triangle-facing half-plane, so that is the half-plane containing~$K$. Intersecting the
three half-planes, $K$ lies in the closed triangle. The tangency equivalence is informal,
standard convex geometry. We have not formalised it, because Mathlib has no notion of an
ellipse or of its tangent lines. The containment argument was informal in version~1. The
Layer-15 file now proves that the Marden ellipse lies in the closed triangle
(\lean{Marden.Classical.l15\_marden\_steiner}). It also proves that if a side line meets the curve
$\{\lean{focalSum}=2A\}$ only at the midpoint, then~(ii) holds for that side
(\lean{Marden.Classical.l15\_strict\_of\_inter}); see Section~\ref{sec:layers}. The formal
\emph{statement} avoids derivatives and angles entirely; the uniqueness \emph{proof}
(Section~\ref{sec:unique}) does use a derivative. Section~\ref{sec:unique} also shows that the
weaker ``supporting line'' condition ($\le$ instead of $<$) already determines the ellipse.

\subsection{The key identity}

\begin{lemma}\label{lem:key}
Let $z_1,z_2$ be the critical points as in Section~\ref{sec:mardenstmt} (so the Vieta relations
hold), and let $m=\frac{a+b}2$. Then $(m-z_1)(m-z_2)=\frac13p'(m)=-\frac{(b-a)^2}{12}$.
\end{lemma}
\begin{proof}
In $p'(m)=(m-b)(m-c)+(m-a)(m-c)+(m-a)(m-b)$, the first two terms cancel because $m-a=-(m-b)$.
The third term is $-(b-a)^2/4$. Formally, the lemma is a single \lean{linear\_combination} of the
two Vieta relations.
\end{proof}

So $(m-z_1)(m-z_2)/(b-a)^2=-1/12$ is a \emph{negative real}. Geometrically, the side direction
$d=b-a$ makes equal angles with $m-z_1$ and $z_2-m$, which is the reflection law of the ellipse at~$m$.
The formal argument turns this into an inequality.

\begin{lemma}\label{lem:reflect}
If $d\ne0$ and $uv=-d^2/12$, then $\nm u+\nm v\le\nm{u+td}+\nm{v+td}$ for every real $t$. The
inequality is strict for $t\ne0$ when $\operatorname{Im}(u/d)\ne0$.
\end{lemma}
\begin{proof}
Write $u=wd$. The hypothesis forces $v=-\mu\bar w d$ with $\mu=1/(12|w|^2)>0$. Hence
$\nm{u+td}+\nm{v+td}=|d|\,(|w+t|+|\mu w-t|)\ge|d|(1+\mu)|w|$ by the triangle inequality, with
equality at $t=0$. When $\operatorname{Im}w\ne0$ and $t\ne0$, the two summands are not on the same ray,
and Mathlib's strict triangle inequality (\lean{norm\_add\_lt\_of\_not\_sameRay}) applies.
\end{proof}

Non-collinearity implies that the focus is off the side line, $\operatorname{Im}\frac{m-z_1}{b-a}\ne0$
(\lean{focus\_off\_side}). This gives~(ii) for the side $ab$. The other sides follow by cyclic
relabelling.

\begin{lemma}\label{lem:focsq}
With $z_1,z_2$ the critical points as in Section~\ref{sec:mardenstmt},
$\bigl(\lean{focalSum}\ z_1\,z_2\,\frac{a+b}2\bigr)^{2}=\frac12|z_1-z_2|^2+\frac19\bigl(|a-b|^2+|b-c|^2+|c-a|^2\bigr)$.
\end{lemma}

The proof combines the parallelogram law with $|m-z_1|\,|m-z_2|=|b-a|^2/12$. The right-hand side
is symmetric in $a,b,c$, so all three midpoints have the same focal sum. This is~(i). It needs only the
Vieta relations, not non-collinearity. Part~(iii) is Vieta.

\subsection{Uniqueness}\label{sec:unique}

\begin{theorem}[uniqueness]\label{thm:unique}
Let $a,b,c,z_1,z_2$ be as in Theorem~\ref{thm:marden}, and let $f_1,f_2\in\C$ and $A'\in\R$ be
arbitrary. Suppose that the three midpoints satisfy $\lean{focalSum}\ f_1\,f_2=2A'$ and that
each side line meets $\{\lean{focalSum}\ f_1\,f_2\le2A'\}$ only at its midpoint. Then
$\{f_1,f_2\}=\{z_1,z_2\}$ and $A'=A$ (\lean{steiner\_unique}). The same conclusion holds if the
strict inequality is weakened to $2A'\le\lean{focalSum}$ along each side (the side is a
supporting line) and the ellipse is assumed non-degenerate, $|f_1-f_2|<2A'$
(\lean{steiner\_inellipse\_unique}).
\end{theorem}

The non-degeneracy premise of the weak form is redundant: it already follows from the three
midpoint equalities and non-collinearity (\lean{nondeg\_of\_midpoints},
Lemma~\ref{lem:mardennonvac}), and \lean{steiner\_unique} discharges it in exactly this way. It
remains an explicit premise of \lean{steiner\_inellipse\_unique} only as a matter of how that
statement is phrased.
So weakening $<$ to $\le$ costs no additional geometric restriction.

\begin{corollary}\label{cor:iff}
For a non-degenerate triangle, a focal ellipse passes through the three midpoints with each side
touching it only there if and only if its foci are the critical points of~$p$ and its major axis is
$2A$ (\lean{steiner\_inellipse\_iff}).
\end{corollary}

The route has three steps.

\begin{lemma}[converse tangency]\label{lem:convtan}
If $d\ne0$, $\nm{u-v}<\nm u+\nm v$, and $t=0$ minimises $t\mapsto\nm{u+td}+\nm{v+td}$, then
$uv=-\lambda d^2$ with $\lambda=\nm u\nm v/\nm d^2\ge0$.
\end{lemma}

The proof sets the derivative at the minimum to zero (\lean{IsLocalMin.hasDerivAt\_eq\_zero}).
A factorisation then leaves two cases. One is the reflection law. The other puts $u$ and $-v$ on
the same ray, which contradicts non-degeneracy.

\begin{lemma}\label{lem:foci}
Let $a,b,c$ be distinct and non-collinear, and let $m_i,d_i$ be their cyclic side midpoints and
directions ($m_1=\frac{a+b}2$, $d_1=b-a$, and cyclically). Suppose the reflection law $(m_i-f_1)(m_i-f_2)=-\lambda_i d_i^2$ holds at all three midpoints, with
$\lambda_i\ge0$, and the three focal sums are equal. Then
$f_1+f_2=\frac23(a+b+c)$, $f_1f_2=\frac13(ab+bc+ca)$ and
$\lambda_1=\lambda_2=\lambda_3=\frac1{12}$ (\lean{foci\_of\_reflection}).
\end{lemma}

Equal focal sums are genuinely needed: the reflection laws with $\lambda_i\ge0$ alone do not
determine the foci. For example, let $a=0$, $b=4$, $c=1+3i$ and let $f_1,f_2$ be the roots of
$z^2-(3+4t+2i)z+(\frac12+10t+4i)$ for real~$t$. Then all three reflection equations hold with
$\lambda_1=\frac3{32}-\frac t8$, $\lambda_2=\frac19-\frac t3$, $\lambda_3=t$, which are
nonnegative for $0\le t\le\frac13$. Only $t=\frac1{12}$ gives equal midpoint focal sums and the
Marden foci (by the lemma, since then $\lambda_i=\frac1{12}$). The reflection identities and the nonnegativity interval are checked exactly in \texttt{code/check\_family.py}; uniqueness of the equal-sum member follows from the preceding lemma. Squaring through
$\nm u\nm v=\lambda\nm d^2$ turns each equality of focal sums into a real linear condition. An
explicit \lean{linear\_combination} certificate shows that it forces $\lambda_i=\lambda_j$
(\lean{lam\_eq\_core}). The certificate multiplies $(\lambda_1-\lambda_2)$ by the square of a quantity
that is nonzero exactly under non-collinearity. Once all the $\lambda_i$ equal $\lambda$, the combination
$\sum_{\mathrm{cyc}}(m_2-m_3)\bigl[(m_1-f_1)(m_1-f_2)+\lambda d_1^2\bigr]$ is independent of
$f_1,f_2$ and equals $(a-b)(a-c)(b-c)(12\lambda-1)/8$ (\lean{lam\_twelfth\_core}). So
$\lambda=\frac1{12}$ is a consequence, not an assumption. The sum and product of the foci then follow
linearly, and $(f_1-z_1)(f_1-z_2)=0$ identifies them.

\paragraph{Scope of uniqueness.} The competitors are all sets
$\{z:|z-f_1|+|z-f_2|=2A'\}$ with arbitrary $f_1,f_2,A'$. Every non-degenerate ellipse in the plane,
including every circle, has this form, and the degenerate cases are excluded automatically
(Lemma~\ref{lem:mardennonvac}). The classical statement that the Steiner inellipse is the unique
non-degenerate \emph{conic} tangent to the sides at their midpoints is not formalised. Stating it would require a
defined notion of non-degenerate conic and of tangency to it, and a bridge showing that a conic of elliptic type has a focal
representation. Mathlib has neither.

\section{Does the statement mean what it says?}\label{sec:correct}

A formal theorem can be true for the wrong reason: its hypotheses may be unsatisfiable, or its
definitions may not mean what the reader thinks. The lemmas below address the points where we
see a risk.

\begin{lemma}[inscribed spheres]\label{lem:inscribed}
Let $X\in\lean{cone V u k}$ with $X\ne V$, and let $ds\ge0$. Then $T=\lean{tangencyPoint V X k d}$ lies on
the sphere $(C_d,ds)$, and $X-T\perp T-C_d$. Moreover no point of the cone lies in the open ball:
$(ds)^2\le\nm{X-C_d}^2$.
\end{lemma}

So an axis sphere with $d>0$ really is inscribed in the cone. It touches every generator and never
crosses the surface. The condition ``tangent to the plane'' in Theorems~\ref{thm:ellipse}
and~\ref{thm:hyperbola} is the distance condition $|\ip{C-p_0}n|=r$. The next lemma connects it
to the pointwise notion.

\begin{lemma}\label{lem:proj}
If $F=C-\ip{C-p_0}n\,n$ and $|\ip{C-p_0}n|=r$, then $F$ lies in the plane, $\nm{F-C}=r$, and
$X-F\perp F-C$ for every $X$ in the plane (\lean{tangent\_of\_proj}, with \lean{proj\_mem\_plane}).
\end{lemma}

\begin{remark}\label{rem:mathlibtangent}
Mathlib has a notion of tangency, \lean{EuclideanGeometry.Sphere.IsTangentAt}, for a sphere and
an affine subspace. \lean{dist\_eq\_of\_isTangentAt} proves the equal-tangents property in that
language. \lean{dandelin\_of\_between\_isTangentAt} is a version of the ellipse theorem that uses
Mathlib spheres and \lean{IsTangentAt}, with betweenness as a hypothesis. The existence theorems
of Section~\ref{sec:dandelin} use the projection form instead. In the main files we have not
restated them with \lean{IsTangentAt}. The Layer-14 file restates them with Mathlib spheres,
\lean{IsTangent} and \lean{IsTangentAt} (Section~\ref{sec:layers}).
\end{remark}

\paragraph{Bridge lemmas.} The module \path{Dandelin/Bridge.lean} ties the definitions to the
classical notions. For $k=\cos\alpha$ with $0\le\alpha\le\pi$ and $X\ne V$,
\lean{mem\_cone\_iff\_angle} says that $X\in\lean{cone V u k}$ iff the generator $VX$ makes angle
$\alpha$ with~$u$, and \lean{mem\_cone2\_iff\_angle} gives angle $\alpha$ or $\pi-\alpha$ for
\lean{cone2}. \lean{generator\_mem\_cone} and \lean{exists\_generator\_of\_mem\_cone} give the
parametric form $V+t\,w$ with $t\ge0$, $\nm w=1$, $\ip wu=k$. \lean{mem\_plane\_iff} is the normal
equation, and \lean{plane\_eq\_mk'} identifies \lean{plane} with Mathlib's affine subspace through
$p_0$ with direction $n^\perp$. \lean{tangencyPoint\_mem\_line} and
\lean{isTangentAt\_tangencyPoint} show that the generator line $VX$ is tangent, in the sense of
\lean{IsTangentAt}, to the axis sphere at \lean{tangencyPoint}. \lean{isTangentAt\_focus₁} and
\lean{isTangentAt\_focus₂} show that the two Dandelin spheres touch the cutting plane, in the same
sense, exactly at \lean{focus₁} and \lean{focus₂}.

\paragraph{Non-vacuity.} The hypotheses of Theorems~\ref{thm:ellipse} and~\ref{thm:hyperbola}
are satisfiable in the intended three-dimensional setting. Non-emptiness is established with the following scopes. Theorem~\ref{thm:regime}(c) proves that the
hyperbola section has points on both nappes (in any~$E$), (d) that the ellipse section is
non-empty when $E$ is finite-dimensional with $\dim E\ge2$, and~(e) the same for the parabola
when $E$ is finite-dimensional with $\dim E\ge3$. For non-circular sections ($m^2<1$) the
directrices are non-empty by Lemma~\ref{lem:infdist}. In version~2 the Layer-14 witness theorems
also prove the hypotheses on explicit configurations in $\R^3$ and $\C$ (Section~\ref{sec:layers}). For Marden:

\begin{lemma}\label{lem:mardennonvac}
For all $a,b,c\in\C$ there are $z_1,z_2$ with $p'=3(X-z_1)(X-z_2)$; this is
\lean{exists\_critical\_points}. If $a,b,c$ are non-collinear and the three midpoints have
common focal sum $2A'$ with respect to some $f_1,f_2$, then $|f_1-f_2|<2A'$
(\lean{nondeg\_of\_midpoints}).
\end{lemma}

The second part applies in particular to $f_i=z_i$, so the Marden ellipse is not degenerate. If it
were a segment or a point, the three midpoints would lie on one line, and then so would $a,b,c$.

\paragraph{Sharpness.} The counterexample of Theorem~\ref{thm:dirconv}(b) is a proved statement.
It shows that the hyperbola converse cannot be stated for a single nappe.

The author has checked each statement in Table~\ref{tab:thms} against its informal counterpart by hand.

\section{Related work}\label{sec:related}

\paragraph{Mathlib.} At tag \texttt{v4.33.1}, \texttt{docs/1000.yaml} lists ``Dandelin's theorem''
and ``Marden's theorem'' with no declaration or author, and the same holds on the master branch as
of 30 September 2026~\cite{Mathlib1000}. A case-insensitive search of the Mathlib sources at that
tag for \texttt{dandelin}, \texttt{marden}, \texttt{inellipse}, \texttt{ellipse} and ``conic
section'' found nothing. Mathlib does contain the Gauss--Lucas theorem
(\texttt{Mathlib/Analysis/Complex/Polynomial/GaussLucas.lean}), which also locates critical
points, and the sphere tangency notion mentioned in Remark~\ref{rem:mathlibtangent}.

\paragraph{Isabelle.} The Archive of Formal Proofs has an entry by Eberl on the geometry of the
critical points of complex polynomials~\cite{EberlAFP}. It covers the Gauss--Lucas and Jensen
theorems. Its bibliography cites Marden's book, but its theory file does not mention Marden's theorem or ellipses. We
searched the titles and abstracts of all 1029 entries in the AFP metadata (mirror repository, commit
\texttt{7de2cdd}, 29 September 2026) for Marden, Dandelin, inellipse, ellipse and conic section, and found no
match. This search did not cover the full text of every theory.

\paragraph{HOL Light.} A search of the HOL~Light repository~\cite{HOLLight} (commit
\texttt{cba9198}, 16 September 2026) for \texttt{marden}, \texttt{dandelin}, \texttt{inellipse}
and \texttt{ellipse} returned no matches. Its \texttt{100/pascal.ml} treats conics, but only
through Pascal's theorem~\cite{BaldPascal}.

\paragraph{Other sources.} The development reports record GitHub searches: code search for
\texttt{Dandelin}, \texttt{Marden}, \texttt{inellipse} and ``Steiner ellipse'' restricted to Lean,
repository search, and a search of Mathlib issues. The only relevant hit was Mathlib issue \#6091,
a wish-list that mentions Marden's theorem. Web searches for formalisations in Lean, Isabelle, Rocq
and HOL found only informal expositions. One search-engine snippet attached to a figure page
on ResearchGate juxtaposes ``the configuration of Marden's theorem'' with the sentence ``We also
provide a machine-checked Lean~4 formalization of the main result''.\footnote{Figure page
\url{https://www.researchgate.net/figure/A-typical-trajectory-of-the-vector-field-in-4-and-the-level-sets-of-the-Lyapunov_fig1_323846042}.
Its title concerns a vector field and a Lyapunov function; fetching it returned HTTP~403. We record it as an unresolved lead, not as evidence of prior work or of its absence.}
The page was not accessible to us, and we could not determine what that formalisation concerns.
% TODO(Josh): check the ResearchGate snippet (figure page 323846042) from a browser; identify the underlying paper.
A search of the Lean Zulip chat for ``Dandelin'' and ``Marden'' on 30 September 2026 returned
no relevant results. We did not search the Mizar library, and Rocq was covered only by web search.
We therefore claim only that we found no earlier complete formalisation of either theorem in the
sources searched.

\section{Verification}\label{sec:verify}

Both libraries were built on a separate machine with Lean \texttt{v4.33.1}, using prebuilt Mathlib
artefacts. Excluding Mathlib, the full Dandelin build takes about 10~s and the Marden build about
17~s, with no warnings. A grep over the library sources (not the scratch files) for \lean{sorry},
\lean{admit}, \lean{native\_decide}, \lean{axiom}, \lean{\#eval}, \lean{run\_cmd},
\lean{set\_option}, \lean{macro}, \lean{elab}, \lean{syntax}, \lean{notation} and
\lean{import Lean} finds nothing. We repeated this grep on the source mirror used for this paper.

Axioms were checked with \lean{\#print axioms} in scratch files outside the libraries, for 63
distinct Dandelin declarations (seven files \texttt{scratch/Ax*.lean}) and 11 Marden declarations.
Each reports only \lean{propext}, \lean{Classical.choice} and \lean{Quot.sound}. The queries cover
every row of Table~\ref{tab:thms}. Five rows (\lean{cone\_tangent\_sq}, \lean{tangent\_of\_proj},
\lean{focalSum\_midpoint\_sq}, \lean{focal\_lt\_of\_mul}, \lean{vieta}) are covered transitively,
because \lean{\#print axioms} follows dependencies and queried theorems use them.
\lean{cone\_outside\_sphere} is used by no other library theorem, so it has its own query
(\texttt{scratch/AxConeOutside.lean}).
The evidence is of two kinds. \emph{Stored logs}: the outputs for 27 of the Dandelin queries
(\texttt{scratch/ax5.out}, \texttt{scratch/ax6.out}, \texttt{scratch/ax\_cone\_outside.out}) and
for all 11 Marden queries (\texttt{ax.log}) are kept as files in the repositories.
\emph{Report-only}: the outputs for the remaining 36 Dandelin queries (the first four phases,
\texttt{scratch/Ax.lean}--\texttt{Ax4.lean}) were recorded only in the development reports; the
query files are kept and can be rerun.
Both repositories also have a \path{gate.sh} that rejects forbidden tokens, builds the library,
generates its own \lean{\#print axioms} queries for every declaration cited in this paper and every
query file in \texttt{scratch/} (85 declarations for Dandelin, 18 for Marden), and fails on any
other axiom. On the published commits both print \texttt{PASS}. In the version-2 commits the
gates also list the Layer-14 and Layer-15 declarations (158 for Dandelin, 56 for Marden). Both
print \texttt{PASS}, and every listed declaration reports only the three standard axioms. This regenerates, on the final
sources, the 36 Dandelin reports that were previously recorded only in the development notes.
\path{tests/tamper.sh} plants six kinds of defect in scratch copies (a \lean{sorry}, a split
\lean{axiom}, a fake \lean{\#eval} report, a macro, a weakened statement, a wrapped extra axiom); in
both repositories the gate rejects all six, each with a nonzero exit status and the expected
failure message. The same holds for the version-2 commits.

Before formalising each statement, we tested it numerically. These checks produced the correction in
Remark~\ref{rem:cos}. The scripts are in \texttt{code/}:
\begin{itemize}
\item \texttt{check\_dandelin.py}, \texttt{check\_phase2.py}, \texttt{check\_phase3.py}: focal sums,
the regime prediction (2000/2000 agreement), the converse, and the focus--directrix relation on 34\,231 points
over 3000 cone--plane pairs (maximum relative error $6.8\cdot10^{-12}$);
\item \texttt{check\_phase4.py}: the hyperbola constant and its converse on 489\,890 points;
\item \texttt{regime\_check.py}, \texttt{degenerate\_check.py}: the constructions of
Theorem~\ref{thm:regime}, the key identity for the directrix converse, the one-nappe counterexample
(69\,106 cases) and the apex sections, all with zero failures;
\item \texttt{check\_marden.py}, \texttt{check\_unique.py}: Lemma~\ref{lem:key} symbolically,
Lemma~\ref{lem:focsq} and the tangency on 2000 random triangles, both \lean{linear\_combination}
certificates of Section~\ref{sec:unique} symbolically, and a sampled perturbation test (a
necessary-condition check for uniqueness, not a verification of it);
\item \texttt{check\_family.py}: exact SymPy check that the one-parameter family of Section~\ref{sec:unique}
satisfies the reflection equations with nonnegative $\lambda_i$ (the focal sums are only sampled).
\end{itemize}
For this paper we re-ran all of the scripts except \texttt{check\_phase4.py}, which takes about nine minutes. The outputs agreed with the recorded
ones. The numerical checks are evidence that the statements are the right ones. They play no role in the proofs.

The step that matters most, comparing the formal statements with the informal ones, cannot be
automated. Sections~\ref{sec:conv} and~\ref{sec:correct} are meant to make that comparison easier.

\subsection{Layers 14 and 15}\label{sec:layers}
\begingroup\sloppy

Version~2 adds a Layer-14 and a Layer-15 library file to each repository:
\path{Dandelin/L14.lean}, \path{Marden/L14.lean}, \path{Dandelin/L15Classical.lean} and
\path{Marden/L15Classical.lean}. The Layer-15 namespaces are \lean{Dandelin.Classical} and
\lean{Marden.Classical}. The gates check selected Layer-14 declarations and the Layer-15 proof
declarations, not every declaration in these files. Every gate-listed declaration reports only the
three standard axioms. The pre-existing proof modules are unchanged; each root module gains two
imports.

\paragraph{Layer 14: statements in Mathlib's language, with witnesses.} Nine main results are
restated with Mathlib-only types and proved from the library theorems through new bridges:
\begin{itemize}
\item seven Dandelin results:
  \begin{itemize}
  \item Theorem~\ref{thm:ellipse}: \lean{Dandelin.l14\_dandelin\_ellipse}, in any real inner
  product space;
  \item Theorem~\ref{thm:hyperbola} together with the double-cone equivalence of
  Theorem~\ref{thm:dirconv}(b): \lean{Dandelin.l14\_dandelin\_hyperbola};
  \item Theorems~\ref{thm:directrix} and~\ref{thm:dirconv}(a) for the ellipse:
  \lean{Dandelin.l14\_directrix\_ellipse};
  \item Theorems~\ref{thm:parabola} and~\ref{thm:dirconv}(a) for the parabola:
  \lean{Dandelin.l14\_parabola};
  \item Theorem~\ref{thm:regime}(b)--(d) in $\R^3$: \lean{Dandelin.l14\_regime\_R3};
  \item Theorem~\ref{thm:apex} in $\R^3$: \lean{Dandelin.l14\_apex\_R3};
  \item Lemma~\ref{lem:inscribed}: \lean{Dandelin.l14\_inscribed};
  \end{itemize}
\item two Marden results: Theorem~\ref{thm:marden} (\lean{Marden.l14\_marden}) and
Theorem~\ref{thm:unique} in its strict form (\lean{Marden.l14\_steiner\_unique}).
\end{itemize}
The Mathlib-only types are as follows.
\begin{itemize}
\item The cone is defined by angles with \lean{InnerProductGeometry.angle}, with the apex
included explicitly, because Mathlib sets the angle of $0$ to $\pi/2$.
\item The plane is \lean{AffineSubspace.mk'} $p_0\,(\R n)^\perp$.
\item Spheres and tangency are \lean{EuclideanGeometry.Sphere}, \lean{IsTangent} and
\lean{IsTangentAt}.
\item The directrix is the intersection of two affine subspaces, and distances to it use
\lean{Metric.infDist}.
\item For Marden we use \lean{¬ Collinear}, \lean{Polynomial.derivative}, \lean{midpoint},
\lean{line[$\R$,x,y]} and \lean{Finset.centroid}.
\end{itemize}
A mechanical check of the elaborated types, with notation turned off, found no constant from
\lean{Dandelin} or \lean{Marden} in any restated statement.

Witness theorems prove every hypothesis on an explicit configuration, together with a
non-degenerate conclusion:
\begin{itemize}
\item \lean{Dandelin.l14w\_dandelin\_ellipse}: $\R^3$ with $\cos\alpha=4/5$; $F_1\ne F_2$, focal
constant $>0$, a section point other than the apex, and non-empty directrices;
\item \lean{Dandelin.l14w\_dandelin\_hyperbola}: $F_1\ne F_2$ and section points on both nappes;
\item \lean{Dandelin.l14w\_parabola}: a section point other than the apex;
\item \lean{Dandelin.l14w\_regime\_inscribed}: a section point other than the apex, at angle
$\alpha$;
\item \lean{Dandelin.l14w\_apex}: the point, line and two-line cases, with distinct non-apex
points on each line;
\item \lean{Marden.l14w\_marden}: the triangle $0,5,1+4i$ with foci $1+2i$ and $3+\frac23i$
(the derivative identity is \lean{Marden.l14w\_derivative}). It proves distinct vertices,
distinct foci, distinct midpoints and a non-degenerate ellipse.
\end{itemize}
Not restated:
\begin{itemize}
\item Corollary~\ref{cor:iff}, which is the conjunction of the two Marden restatements;
\item the weak form of Theorem~\ref{thm:unique};
\item Theorem~\ref{thm:regime}(a) and~(e), and the versions in general dimension;
\item the one-nappe counterexample of Theorem~\ref{thm:dirconv}(b);
\item the internal lemmas. Lemma~\ref{lem:proj} has a Mathlib-only form as the bridge
\lean{Dandelin.l14\_isTangent\_plane\_iff}.
\end{itemize}
No restatement turned out false or needed an extra hypothesis.

\paragraph{Layer 15: comparison with textbook statements.} Another AI agent collected the
statements of the classical theorems from the literature, with citations: Dandelin, Hilbert and
Cohn-Vossen, Kalman, Marden, Siebeck, and standard reference pages. An AI model was given only
these textbook statements. It saw no code and no part of this paper, and its transcript was
audited. It wrote 13 Lean statements with Mathlib notions and its own definitions: 12 theorems and
one well-formedness obligation, namely that the directrix is a line. Its Dandelin sphere is
\emph{any} sphere of positive radius meeting every generator in exactly one point and the plane
in exactly one point. It marked three ambiguities and no extra hypotheses.

A comparer copied these statements verbatim. For each one it either proved it from the development
or classified it. It added no hypotheses and left the library unchanged. The outcome: 11 covered,
1 narrower, 1 out of scope, and none false as written.
\begin{itemize}
\item \emph{Covered and proved} (11): \lean{cov\_Conics1\_a}, \lean{cov\_Conics1\_b},
\lean{cov\_Conics2\_a}, \lean{cov\_Conics2\_b}, \lean{cov\_Conics3\_obligation} (the obligation is
proved, not assumed), \lean{cov\_Conics3\_a}, \lean{cov\_Conics3\_b} and
\lean{cov\_Conics3\_parabola\_unique} in \lean{Dandelin.Classical}; \lean{cov\_Conics4\_a},
\lean{cov\_Conics5} and \lean{cov\_Conics5\_double\_root} in \lean{Marden.Classical}.
\item \emph{Narrower} (1): for \lean{classical\_Conics4\_b}, only the centre clause is proved
(\lean{Marden.Classical.cov\_Conics4\_b\_part}: the inellipse is centrally symmetric about
$(a+b+c)/3$). The area formula $\pi/(3\sqrt3)\cdot\operatorname{area}(abc)$ and the strict maximality
of the area among inscribed ellipses are missing. This paper does not claim them, and the libraries
contain no measure theory.
\item \emph{Out of scope} (1): \lean{classical\_Conics3\_c} is not proved. This is the planar
converse: in an abstract plane, the locus $|XF|=e\cdot\mathrm{dist}(X,d)$ is an ellipse, parabola
or hyperbola in standard form, with focus $F$ and directrix $d$. Theorem~\ref{thm:dirconv} is a
converse on the cone, not this statement. The comparer checked by hand that the statement as
written is true.
\item \emph{Writer statement false as written}: none.
\end{itemize}
New proofs, which are not in the main files:
\begin{itemize}
\item \lean{Dandelin.Classical.l15\_tangent\_cone\_iff}: in $\R^3$, a sphere of positive radius
meets every generator in exactly one point if and only if it is an axis sphere $(C_d,|d|s)$ with
$d\ne0$. Hence the uniqueness in Conics1\_a, Conics2\_a and Conics3\_parabola\_unique holds among
all such spheres, not only among axis spheres.
\item \lean{Dandelin.Classical.l15\_gen\_segment}: the segment of every generator between the two
contact circles has length $|d_1-d_2|\cos\alpha$, which equals the focal constant.
\item \lean{Dandelin.Classical.l15\_isGreatest}: the focal constant of the ellipse is the greatest
distance between two of its points, that is, the major axis.
\item The transverse axis, in \lean{Dandelin.Classical.cov\_Conics2\_b}: the focal constant of the
hyperbola is the least distance between its two branches.
\item \lean{Dandelin.Classical.l15\_span\_circle}, \lean{Dandelin.Classical.l15\_directrix\_eq}:
the affine span of a contact circle is the plane $\sigma_d$, and its intersection with the cutting
plane is the directrix $\ell_d$, which is a line when $m^2<1$ (equivalently, $\beta<\pi/2$).
\item \lean{Dandelin.Classical.l15\_planeAxisAngle}: $\cos\beta=\sqrt{1-m^2}$, so the paper's $e$ is
$\cos\beta/\cos\alpha$.
\item \lean{Marden.Classical.l15\_strict\_of\_inter}: if a side line meets the ellipse curve only at
the midpoint, then the strict tangency~(ii) of Theorem~\ref{thm:marden} holds.
\item \lean{Marden.Classical.l15\_marden\_steiner}: the Marden ellipse lies in the closed
triangle.
\item \lean{Marden.Classical.l15\_equilateral\_iff}: $p'$ has a double root if and only if the
triangle is equilateral.
\end{itemize}
The classical definitions and statements are reproduced verbatim in the Layer-15 files. The final
gates list 158 declarations for Dandelin (85 published, 17 added by Layer~14, 56 by Layer~15) and
56 for Marden (18, 6 and 32). Both pass, and the tampering suites reject every planted defect.
\par\endgroup

\section{Limitations and further work}

\begin{itemize}
\item \emph{Tangency.} Neither development uses a general notion of a tangent line or plane to a
conic. For Dandelin, tangency of spheres and planes is expressed by distance and projection
(Lemma~\ref{lem:proj}). In the main files it is connected to Mathlib's \lean{IsTangentAt} only in the
normalised version (Remark~\ref{rem:mathlibtangent}). The Layer-14 restatements use \lean{IsTangent}
and \lean{IsTangentAt} throughout (Section~\ref{sec:layers}). For Marden, tangency is metric
(Section~\ref{sec:metrictan}). The equivalence with the differential notion is argued but not
formalised.
\item \emph{Uniqueness among conics.} Theorem~\ref{thm:unique} is uniqueness among focal
ellipses, not among all conics.
\item \emph{Not formalised.} We do not formalise the area of the Steiner inellipse
($\pi/(3\sqrt3)$ times the area of the triangle) or its maximality among inscribed ellipses.
We also do not formalise the planar converse that the locus $|XF|=e\cdot\mathrm{dist}(X,d)$ is a
conic in standard form (Section~\ref{sec:layers}).
\item \emph{Cones.} Only right circular cones are treated. The regime equivalences that need
finite dimension and $\dim E\ge3$ are proved under those hypotheses, and the two-line apex section only in
$\dim E=3$. For planes through the apex, only the double cone is treated; the one-nappe versions,
which give rays, are not.
\item \emph{Reference geometry.} Mathlib has no ellipse, parabola or hyperbola, so the
conics appear only as loci given by explicit equations. A Mathlib contribution could take the
equation form as it stands, but would benefit from agreed definitions of these curves and bridges
to them.
\end{itemize}

\section*{Acknowledgements and use of AI}
The formal proofs, the numerical checking scripts and a draft of this paper were prepared with the
help of large-language-model AI assistants. Every proof is checked by the Lean kernel, and the
author takes responsibility for all content. The author compared each formal statement in
Table~\ref{tab:thms} with its informal counterpart by hand. AI agents carried out the Layer-14 and
Layer-15 checks of version~2 (Section~\ref{sec:layers}). This includes collecting the classical
statements from the literature, which no human has yet checked against the sources.

\begin{thebibliography}{99}
\bibitem{BaldPascal} J.~Bald, Pascal's hexagon theorem over an arbitrary field, formalised in
Lean~4, preprint, 2026.
\bibitem{Dandelin} G.~P.~Dandelin, M\'emoire sur quelques propri\'et\'es remarquables de la
focale parabolique, \emph{Nouveaux m\'emoires de l'Acad\'emie royale des sciences et
belles-lettres de Bruxelles} 2 (1822), 170--202.
\bibitem{EberlAFP} M.~Eberl, Two theorems about the geometry of the critical points of a complex
polynomial, \emph{Archive of Formal Proofs}, November 2023,
\url{https://isa-afp.org/entries/Polynomial_Crit_Geometry.html}.
\bibitem{HilbertCV} D.~Hilbert and S.~Cohn-Vossen, \emph{Geometry and the Imagination}, translated
by P.~Nemenyi, Chelsea, New York, 1952.
\bibitem{HOLLight} J.~Harrison, HOL~Light repository, \url{https://github.com/jrh13/hol-light},
accessed 30 September 2026.
\bibitem{Kalman} D.~Kalman, An elementary proof of Marden's theorem, \emph{Amer. Math. Monthly}
115 (2008), no.~4, 330--338.
\bibitem{Marden1945} M.~Marden, A note on the zeros of the sections of a partial fraction,
\emph{Bull. Amer. Math. Soc.} 51 (1945), no.~12, 935--940.
\bibitem{MardenBook} M.~Marden, \emph{Geometry of Polynomials}, 2nd ed., Mathematical Surveys~3,
American Mathematical Society, 1966.
\bibitem{Mathlib1000} The mathlib Community, \texttt{docs/1000.yaml}, Mathlib repository,
\url{https://github.com/leanprover-community/mathlib4}, accessed 30 September 2026.
\bibitem{MathlibCPP} The mathlib Community, The Lean mathematical library, in \emph{Proc.\ 9th
ACM SIGPLAN Int.\ Conf.\ on Certified Programs and Proofs (CPP 2020)}, 2020, 367--381.
\bibitem{MindaPhelps} D.~Minda and S.~Phelps, Triangles, ellipses, and cubic polynomials,
\emph{Amer. Math. Monthly} 115 (2008), no.~8, 679--689.
\bibitem{Lean4} L.~de~Moura and S.~Ullrich, The Lean~4 theorem prover and programming language,
in \emph{Automated Deduction -- CADE~28}, Lecture Notes in Computer Science, Springer, 2021,
625--635.
\bibitem{Siebeck} Siebeck, Ueber eine neue analytische Behandlungsweise der Brennpunkte,
\emph{J. Reine Angew. Math.} 64 (1865), 175--182.
\end{thebibliography}

\end{document}
